\section{Extensión triádica de Euler y conservación de formas}
\label{eul:capitulo}

La representación cuadrática del TRIT adquiere contenido operatorio al
identificar los generadores seleccionados por APP y por el calendario del
TPK. Este capítulo reconstruye esos antecedentes antes de expresar sus
realizaciones exponenciales. El orden distingue la operación discreta,
su representación matricial y la función analítica que permite describir
su evolución en un parámetro real.

Tres construcciones intervienen: el ciclo multiplicativo de orden tres
en APP, el transporte elemental que conserva una coordenada de
prolongación y la incidencia entre modos que desciende del calendario
trítico. Sus generadores satisfacen, respectivamente, las relaciones
elíptica, parabólica e hiperbólica. La forma alternante del plano permite
después establecer una conservación común, con las diferencias métricas
y de orientación de cada realización expresamente conservadas.

\subsection{Ciclo multiplicativo de APP y generador elíptico}
\label{eul:coxeter}

En \(\mathbb Z/9\mathbb Z\), la multiplicación por cuatro es invertible
y tiene orden tres:
\[
 4^2=7\pmod9,\qquad 4^3=1\pmod9.
\]
Sobre las unidades aparecen los ciclos \((1,4,7)\) y \((2,8,5)\).
Las clases \(0,3,6\) quedan fijas. La selección del multiplicador y sus
órbitas pertenece, por tanto, a la aritmética nonádica de APP.

\subsubsection{Acción contragrediente sobre las dos evaluaciones}

Sean \(\bar S(x,y)=x+y\) y \(\bar P(x,y)=xy\), con valores en
\(\mathbb Z/9\mathbb Z\). Para \(a=4\),
\begin{equation}
 \bar S(ax,ay)=a\bar S(x,y),\qquad
 \bar P(ax,ay)=a^2\bar P(x,y)=a^{-1}\bar P(x,y).
\label{eul:eq:contragrediente}
\end{equation}
La primera identidad procede de la distributividad; la segunda, de
\(a^3=1\). La suma y el producto recorren orientaciones opuestas del
mismo ciclo residual. En las evaluaciones enteras, los cocientes
nonádicos permanecen asociados al levantamiento de estas clases.

\subsubsection{Módulo de aumento y matriz del ciclo}

Fijado uno de los ciclos \((r,4r,7r)\), considérese el módulo libre
generado por \(e_r,e_{4r},e_{7r}\). Su submódulo de aumento es
\[
 L=\{x_0e_r+x_1e_{4r}+x_2e_{7r}:x_0+x_1+x_2=0\}.
\]
Una base es \(b_1=e_r-e_{4r}\), \(b_2=e_{4r}-e_{7r}\). El ciclo
manda \(b_1\) a \(b_2\) y \(b_2\) a \(-b_1-b_2\). Su matriz,
actuando sobre vectores columna, es
\[
 C=\begin{pmatrix}0&-1\\1&-1\end{pmatrix}.
\]
El producto escalar heredado del módulo de tres coordenadas tiene matriz
de Gram
\[
 G_{\mathrm e}=\begin{pmatrix}2&-1\\-1&2\end{pmatrix}.
\]
Se obtiene así el retículo \(A_2\), con su forma determinada por las
diferencias de posiciones del ciclo.

\begin{proposicion}[Normalización elíptica del ciclo APP]
\label{eul:prop:eliptico}
El ciclo satisface
\[
 C^2+C+I=0,\qquad C^\mathsf TG_{\mathrm e}C=G_{\mathrm e}.
\]
En \(L\otimes_{\mathbb Z}\mathbb R\), el operador
\[
 J_{+1}=\frac{2C+I}{\sqrt3}
\]
tiene traza cero y satisface \(J_{+1}^2=-I\).
\end{proposicion}
\begin{proof}
La multiplicación da
\(C^2=\bigl(\begin{smallmatrix}-1&1\\-1&0\end{smallmatrix}\bigr)\),
de donde resulta \(C^2+C+I=0\). La permutación cíclica conserva el
producto escalar de las tres coordenadas y su submódulo de aumento;
en la base elegida esto equivale a la identidad de Gram. Finalmente,
\((2C+I)^2=4C^2+4C+I=-3I\), y
\(\operatorname{tr}(2C+I)=0\).
\end{proof}

La normalización real proporciona una raíz operatoria de \(-I\) a
partir del ciclo ya construido. La identidad puede escribirse también
\[
 C=-\frac12I+\frac{\sqrt3}{2}J_{+1}.
\]
Su interpretación angular se obtiene después de fijar la realización
analítica del tipo elíptico.

\subsection{Transporte de umbral y generador parabólico}
\label{eul:parabolico}

Considérense las coordenadas enteras \((x,c)\), donde \(c\) se conserva
como incremento registrado. El transporte elemental
\[
 (x,c)\longmapsto(x+c,c)
\]
se representa mediante
\[
 U_0=I+N,
 \qquad N=\begin{pmatrix}0&1\\0&0\end{pmatrix}.
\]
Éste es el modelo lineal del transporte de umbral utilizado en la
realización cuadrática. La división balanceada y su normalización,
desarrolladas en el capítulo~\ref{trit:capitulo}, conservan sus reglas
propias cuando se aplican a una evaluación entera.

\begin{proposicion}[Iteración y reversibilidad del transporte]
\label{eul:prop:parabolico}
Se cumple \(N^2=0\), \(N\ne0\) y, para todo \(m\in\mathbb Z\),
\[
 U_0^m=I+mN,
 \qquad U_0^m(x,c)=(x+mc,c).
\]
El transporte conserva el determinante y admite inversa \(I-N\).
\end{proposicion}
\begin{proof}
La imagen de \(N\) está contenida en la primera coordenada, sobre la
que \(N\) se anula; por ello \(N^2=0\). La identidad
\((I+aN)(I+bN)=I+(a+b)N\) prueba las potencias positivas, la inversa y
las potencias negativas. El determinante de \(I+mN\) es uno.
\end{proof}

El índice parabólico \(\tau=0\) corresponde a un generador de cuadrado
nulo. El operador \(N\) transporta un dato lineal y puede ser distinto
de cero sobre un estado concreto. La firma nula, el operador nulo y un
generador nilpotente son tres objetos distintos.

\subsection{Incidencia de modos derivada del calendario TPK}
\label{eul:incidencia-modos}

La construcción de la sección~\ref{sec:cal-incidencia} proporciona
la matriz~\eqref{eq:cal-incidencia}. Se explicita aquí su acción
primitiva para enlazar el conteo trítico con la representación cuadrática
hiperbólica que se desarrolla a continuación.

El selector trítico de fase toma, sobre las posiciones \(1,\ldots,9\),
los valores \((+1,-1,0)\) repetidos tres veces. Sobre el modo operativo
se utiliza la regla
\[
 +1:\ m\mapsto\Sigma,\qquad
 -1:\ m\mapsto\Pi,\qquad
 0:\ m\mapsto m.
\]
El bloque primitivo produce, después de cada actualización, la sucesión
de modos \((\Sigma,\Pi,\Pi)\). El transporte espacial y la actualización
del registro permanecen en la transición completa del TPK; aquí se lee
su incidencia modal.

\begin{teorema}[Incidencia primitiva del calendario]
\label{eul:teo:fav}
Al contar pares consecutivos con cierre cíclico en
\((\Sigma,\Pi,\Pi)\), la matriz obtenida en el orden \((\Sigma,\Pi)\)
es
\[
 F_{\rm av}=\begin{pmatrix}0&1\\1&1\end{pmatrix}.
\]
En calendarios de longitud \(9,27,108\), los conteos son
\(3F_{\rm av},9F_{\rm av},36F_{\rm av}\), respectivamente.
\end{teorema}
\begin{proof}
Los tres pares son \(\Sigma\to\Pi\), \(\Pi\to\Pi\) y
\(\Pi\to\Sigma\). Cada uno aparece una vez. El par
\(\Sigma\to\Sigma\) aparece cero veces. Esto determina las cuatro
entradas. La repetición del calendario multiplica todos los conteos
por el número de bloques primitivos, que es \(3,9,36\) en los casos
indicados.
\end{proof}

El lector de hojas asocia \(\Sigma\) a la lectura areal y \(\Pi\) a
la volumétrica. Esta asignación transporta la incidencia obtenida a
los dos tipos de lectura. La matriz resulta del calendario; su análisis
espectral se efectúa una vez construidas sus entradas.

Los conteos de calendarios repetidos y las potencias matriciales tienen
funciones diferentes. \(36F_{\rm av}\) cuenta los pares recorridos
durante un calendario de \(108\) eventos. En cambio,
\(F_{\rm av}^{108}\) cuenta cadenas admisibles de longitud \(108\)
en el grafo de incidencia. El primer objeto sigue una cronología;
el segundo pertenece a su envolvente de adyacencias.

\subsubsection{Realización hiperbólica de la incidencia}

Dos composiciones de incidencia producen
\[
 A=F_{\rm av}^2=\begin{pmatrix}1&1\\1&2\end{pmatrix}.
\]
Sus coeficientes cuentan recorridos de dos aristas. Se tiene
\(\det A=1\), \(\operatorname{tr}A=3\) y, por multiplicación directa,
\(A^2-3A+I=0\).

\begin{proposicion}[Generador hiperbólico y forma invariante]
\label{eul:prop:hiperbolico}
El operador
\[
 J_{-1}=\frac{2A-3I}{\sqrt5}
\]
tiene traza cero y cuadrado \(I\). La forma simétrica
\[
 G_{\mathrm h}=\begin{pmatrix}2&1\\1&-2\end{pmatrix}
\]
tiene firma \((1,1)\) y satisface
\(A^\mathsf TG_{\mathrm h}A=G_{\mathrm h}\).
\end{proposicion}
\begin{proof}
De \(A^2-3A+I=0\) se obtiene
\((2A-3I)^2=4A^2-12A+9I=5I\). La traza del numerador es cero.
Además,
\[
 G_{\mathrm h}A=\begin{pmatrix}3&4\\-1&-3\end{pmatrix},
 \qquad
 A^\mathsf TG_{\mathrm h}A
 =\begin{pmatrix}2&1\\1&-2\end{pmatrix}.
\]
El determinante de \(G_{\mathrm h}\) es \(-5\), por lo que sus dos
autovalores tienen signos opuestos.
\end{proof}

La realización es lorentziana en dimensión \(1+1\): queda especificado
un plano real, una forma de firma \((1,1)\) y un transporte que la
conserva. La realización física de un espacio-tiempo utiliza después
sus campos, dimensiones y mapas de transducción correspondientes.

\subsection{Identidad exponencial de los tres regímenes}
\label{eul:identidad}

Los tres operadores concretos actúan en espacios de procedencia distinta:
el módulo de aumento de un ciclo APP, el plano del transporte de umbral
y el plano de incidencia modal. Sus normalizaciones satisfacen
\[
 J_\tau^2=-\tau I,\qquad \operatorname{tr}J_\tau=0,
 \qquad \tau\in\{+1,0,-1\}.
\]
Esta relación determina sus potencias y permite construir las funciones
de evolución por series convergentes.

\begin{definicion}[Funciones del régimen cuadrático]
Para \(t\in\mathbb R\), se definen
\[
 c_\tau(t)=\sum_{k=0}^{\infty}
       \frac{(-\tau)^kt^{2k}}{(2k)!},\qquad
 s_\tau(t)=\sum_{k=0}^{\infty}
       \frac{(-\tau)^kt^{2k+1}}{(2k+1)!}.
\]
Estas funciones realizan analíticamente el tipo cuadrático ya obtenido.
\end{definicion}

\begin{teorema}[Extensión triádica de la identidad de Euler]
\label{eul:teo:euler}
Para los generadores anteriores,
\begin{equation}
 E_\tau(t):=\exp(tJ_\tau)
 =c_\tau(t)I+s_\tau(t)J_\tau.
\label{eul:eq:euler}
\end{equation}
Las especializaciones de \((c_\tau,s_\tau)\) son
\((\cos,\sin)\), \((1,t)\) y \((\cosh,\sinh)\), en el orden
\(\tau=+1,0,-1\). Además,
\[
 c_\tau(t)^2+\tau s_\tau(t)^2=1,
 \qquad E_\tau(t)^{-1}=E_\tau(-t).
\]
\end{teorema}
\begin{proof}
La serie exponencial converge absolutamente en norma matricial.
La relación cuadrática da
\(J_\tau^{2k}=(-\tau)^kI\) y
\(J_\tau^{2k+1}=(-\tau)^kJ_\tau\). Separar los términos pares e
impares produce \eqref{eul:eq:euler}. Las tres especializaciones se
reconocen en sus series. Como las matrices \(tJ_\tau\) y \(-tJ_\tau\)
conmutan, sus exponenciales tienen producto \(I\). La paridad de las
series da
\[
 E_\tau(t)E_\tau(-t)
 =\bigl(c_\tau(t)^2+\tau s_\tau(t)^2\bigr)I,
\]
lo que prueba la identidad escalar y la inversa.
\end{proof}

Los números que aparecen al evaluar funciones analíticas tienen aquí
estatuto de reconocimiento de operadores previamente construidos.
La generación regional coinductiva conserva su emisor, sus elevaciones
y sus cilindros. En particular, estas series describen la realización
exponencial; la selección de los prefijos regionales se realiza mediante
el TPK.

\subsubsection{Ley de composición e invariantes cuadráticos}

La conmutación de los múltiplos de un mismo generador da
\(E_\tau(t)E_\tau(s)=E_\tau(t+s)\). Comparando coeficientes de
\(I,J_\tau\),
\begin{align}
 c_\tau(t+s)&=c_\tau(t)c_\tau(s)-\tau s_\tau(t)s_\tau(s),\\
 s_\tau(t+s)&=s_\tau(t)c_\tau(s)+c_\tau(t)s_\tau(s).
\end{align}
La conjugación cuadrática cambia \(t\) por \(-t\) dentro del régimen
fijado. La inversión orientada de una sucesión trítica, por su parte,
puede cambiar los índices de régimen y debe transportarse mediante la
involución correspondiente del TPK.

En el tipo elíptico, la identidad conservada es \(a^2+b^2=1\).
En el hiperbólico, \(a^2-b^2=1\), y las coordenadas \(u=a+b\),
\(v=a-b\) verifican \(uv=1\). La evolución hiperbólica tiene
\((u,v)=(e^t,e^{-t})\): la expansión de una coordenada se acompaña
de la contracción de la otra. En el tipo parabólico,
\(E_0(t)=I+tN\), con transporte lineal y determinante uno.

\subsubsection{Evaluaciones sobre los operadores discretos}

La representación del ciclo elíptico da
\[
 C=E_{+1}(2\pi/3),\qquad
 C^n=E_{+1}(2\pi n/3),\quad n\in\mathbb Z.
\]
De manera particular,
\[
 E_{+1}(\pi)+I=0,\qquad E_{+1}(2\pi)=I,
 \qquad E_0(1)=I+N.
\]
La coordenada angular \(\pi\) expresa analíticamente la media vuelta;
su identificación con la salida regional HMT corresponde al lector
coinductivo desarrollado en su lugar causal.

La incidencia \(F_{\rm av}\) tiene polinomio característico
\(x^2-x-1\). Su raíz positiva \(\lambda\) satisface
\(\lambda^2=\lambda+1\); al reconocerla como razón áurea se escribe
\(\lambda=\varphi\). Con \(\eta=2\log\lambda\),
\[
 \cosh\eta=\frac32,\qquad \sinh\eta=\frac{\sqrt5}{2},
 \qquad E_{-1}(\eta)=A=F_{\rm av}^2.
\]
La comprobación usa \(\lambda^{-1}=\lambda-1\), de donde
\(\lambda^2+\lambda^{-2}=3\) y
\(\lambda^2-\lambda^{-2}=\sqrt5\). La matriz se obtuvo antes por
conteo de transiciones; el reconocimiento espectral caracteriza su
escala. La igualdad con la coordenada regional generada se conserva
con el lector y la demostración de dicha región.

Asimismo, la identidad escalar \(\exp(I)=eI\) caracteriza la evaluación
unitaria de la exponencial. Su función en la representación analítica
es distinta de la construcción coinductiva de la región correspondiente.

\subsection{Resolución polinómica de los regímenes}
\label{eul:proyectores}

En la suma directa de los tres espacios reales, defínanse
\[
 \mathbb J=J_{+1}\oplus N\oplus J_{-1},\qquad Q=\mathbb J^2.
\]
En esos sumandos, \(Q\) actúa como \(-I,0,I\), respectivamente.
Por tanto \(Q^3=Q\) y \(\mathbb J^6=\mathbb J^2\).

\begin{proposicion}[Idempotentes de régimen]
\label{eul:prop:proyectores}
Los operadores
\[
 p_{\mathrm e}=\frac{Q^2-Q}{2},\qquad
 p_{\mathrm p}=I-Q^2,\qquad
 p_{\mathrm h}=\frac{Q^2+Q}{2}
\]
son idempotentes ortogonales en sentido algebraico y suman \(I\).
Sus imágenes son los tres sumandos de la representación.
\end{proposicion}
\begin{proof}
Los tres polinomios toman los valores \((1,0,0)\), \((0,1,0)\) y
\((0,0,1)\) sobre \(-1,0,+1\). Sus cuadrados menos ellos mismos y
sus productos cruzados se anulan en esos tres puntos, por lo que son
múltiplos de \(X^3-X\). Sustituir \(Q\) demuestra idempotencia y
anulación de los productos cruzados. La suma de los tres polinomios
es uno. Su acción en cada bloque identifica las imágenes.
\end{proof}

La exponencial completa queda expresada por
\begin{align}
 \exp(t\mathbb J)={}&p_{\mathrm e}(\cos t\,I+\sin t\,\mathbb J)
 +p_{\mathrm p}(I+t\mathbb J)\notag\\
 &+p_{\mathrm h}(\cosh t\,I+\sinh t\,\mathbb J).
\end{align}
Esta resolución distingue exactamente los regímenes. Las transiciones
entre ellos pertenecen a los mapas del TPK sobre estados con orientación,
modo y memoria; la suma directa conserva sus dominios separados.

\subsection{Compresión nonádica y registro complementario}
\label{eul:compresion}

La conservación entre una lectura y su registro complementario procede
de la estructura de nueve fases. Sea \(V\) un espacio con producto
escalar y \(C:V\to V\) un transporte lineal. En \(V^9\), defínanse
\[
 S_C(v_0,\ldots,v_8)=(Cv_8,v_0,\ldots,v_7),
 \qquad Jv=\frac13(v,\ldots,v).
\]
La inclusión \(J\) es isométrica porque reúne nueve copias con factor
\(1/3\). Su adjunto suma las componentes y divide por tres. La
compresión sobre esa inclusión es
\[
 T=J^*S_CJ=\frac{8I+C}{9}.
\]
Ocho componentes conservan \(v\) y una atraviesa el transporte \(C\).
El complemento ortogonal de la inclusión registra
\[
 (I-JJ^*)S_CJv
 =\frac1{27}\bigl(8(C-I)v,-(C-I)v,\ldots,-(C-I)v\bigr).
\]
El cuadrado de su norma es
\(8\|(C-I)v\|^2/81\). El mismo contenido vectorial puede representarse
isométricamente por
\[
 D=\frac{\sqrt8}{9}(C-I),
\]
manteniendo el registro de nueve componentes cuando se utilizan sus
etiquetas individuales.

\begin{proposicion}[Recuperación algebraica desde dos lecturas]
\label{eul:prop:recuperacion}
Para todo \(v\in V\),
\begin{equation}
 v=Tv-\frac1{\sqrt8}Dv.
\label{eul:eq:recuperacion}
\end{equation}
Si \(C\) es isométrico, se cumple además
\[
 \|v\|^2=\|Tv\|^2+\|Dv\|^2.
\]
\end{proposicion}
\begin{proof}
La primera identidad resulta de
\((8I+C-(C-I))/9=I\). Para la segunda, al expandir los productos
se obtiene
\[
 T^*T+D^*D=\frac{72I+9C^*C}{81}.
\]
La hipótesis \(C^*C=I\) da la identidad. Los términos mixtos se
cancelan con coeficientes \(+8\) y \(-8\).
\end{proof}

La recuperación utiliza los registros vectoriales, incluidas sus
correlaciones. Las intensidades \(\|Tv\|^2\) y \(\|Dv\|^2\)
proporcionan el balance escalar, mientras que la fórmula inversa utiliza
\(Tv\) y \(Dv\) completos.

\subsection{Conservación del área orientada}
\label{eul:area}

En cualquier plano real orientado, sea
\[
 \Omega=\begin{pmatrix}0&1\\-1&0\end{pmatrix}.
\]
Para una matriz \(L=\bigl(\begin{smallmatrix}a&b\\c&d\end{smallmatrix}\bigr)\),
la multiplicación directa da
\begin{equation}
 L^\mathsf T\Omega L=(ad-bc)\Omega=(\det L)\Omega.
\label{eul:eq:determinante-area}
\end{equation}
Como los generadores tríticos tienen traza cero,
\(\det E_\tau(t)=1\). Sus realizaciones conservan, por tanto, la forma
alternante, aunque sus formas métricas sean diferentes.

\begin{teorema}[Conservación nonádica común a los tres regímenes]
\label{eul:teo:area}
Para \(C=E_\tau(t)\), sean
\(T=(8I+C)/9\) y \(D=\sqrt8(C-I)/9\). Entonces
\begin{equation}
 \Omega=T^\mathsf T\Omega T+D^\mathsf T\Omega D.
\label{eul:eq:balance-area}
\end{equation}
\end{teorema}
\begin{proof}
La expansión da
\[
 T^\mathsf T\Omega T+D^\mathsf T\Omega D
 =\frac{72\Omega+9C^\mathsf T\Omega C}{81}.
\]
Los términos \(C^\mathsf T\Omega\) y \(\Omega C\) se cancelan.
Por \eqref{eul:eq:determinante-area},
\(C^\mathsf T\Omega C=\Omega\), y el resultado es \(\Omega\).
\end{proof}

Con \(c=c_\tau(t)\), las contribuciones determinantes son
\begin{equation}
 \det T=\frac{65+16c}{81},\qquad
 \det D=\frac{16(1-c)}{81},\qquad
 \det T+\det D=1.
\label{eul:eq:determinantes}
\end{equation}
En efecto, para \(\det C=1\),
\(\det(aI+C)=a^2+a\operatorname{tr}C+1\), y
\(\operatorname{tr}C=2c\).

\begin{tablacompleta}
\caption{Contribuciones al área orientada de tres transportes generados.}
\begin{tabular}{lccc}
\toprule
Transporte&\(\operatorname{tr}C\)&\(\det T\)&\(\det D\)\\
\midrule
Ciclo APP \(C\)&$-1$&$19/27$&$8/27$\\
Umbral \(I+N\)&$2$&$1$&$0$\\
Incidencia \(F_{\rm av}^2\)&$3$&$89/81$&$-8/81$\\
\bottomrule
\end{tabular}
\notatabla{Cada par suma uno. El signo registra orientación; la tabla
se refiere a área alternante, no a probabilidades.}
\end{tablacompleta}

En el transporte parabólico, \(D\) puede ser no nulo y tener rango uno:
su área es cero y su registro contiene información lineal. En el
hiperbólico, el complemento tiene orientación opuesta. En el elíptico,
la forma positiva propia del ciclo permite, además, el balance de norma.
Las tres realizaciones comparten la identidad alternante con sus tipos
conservados.

\subsection{Formas variables y restitución entre planos sucesivos}
\label{eul:formas-variables}

Considérese una sucesión de planos \(V_n\) con matrices simétricas
positivas definidas \(G_n\). Sean
\(B_n,C_n:V_n\to V_{n+1}\) dos transportes que satisfacen
\[
 B_n^\mathsf TG_{n+1}B_n=G_n,
 \qquad C_n^\mathsf TG_{n+1}C_n=G_n.
\]
El transporte \(B_n\) identifica las coordenadas entre los dos planos;
\(C_n\) incluye la evolución adicional. Defínanse
\[
 T_n=\frac{8B_n+C_n}{9},\qquad
 D_n=\frac{\sqrt8}{9}(C_n-B_n).
\]

\begin{teorema}[Conservación y recuperación con forma variable]
\label{eul:teo:forma-variable}
Se cumplen
\begin{align}
 G_n&=T_n^\mathsf TG_{n+1}T_n+D_n^\mathsf TG_{n+1}D_n,
 \label{eul:eq:balance-variable}\\
 v&=B_n^{-1}\left(T_nv-\frac1{\sqrt8}D_nv\right).
 \label{eul:eq:recuperacion-variable}
\end{align}
\end{teorema}
\begin{proof}
En la expansión de la primera expresión, los términos cruzados se
cancelan. Quedan \(72B_n^\mathsf TG_{n+1}B_n/81\) y
\(9C_n^\mathsf TG_{n+1}C_n/81\), cuya suma es \(G_n\).
La identidad \(T_n-D_n/\sqrt8=B_n\) prueba la segunda fórmula.
La positividad de las formas hace inyectivo a \(B_n\); entre planos
de igual dimensión, es invertible.
\end{proof}

Para \(I_n(v)=\tfrac12v^\mathsf TG_nv\),
\(v_{n+1}=T_nv_n\) y \(m_n=D_nv_n\), se obtiene
\[
 I_n(v_n)=I_{n+1}(v_{n+1})+I_{n+1}(m_n).
\]
La suma de \(N\) pasos da
\begin{equation}
 I_0(v_0)=I_N(v_N)+\sum_{n=0}^{N-1}I_{n+1}(m_n).
\label{eul:eq:telescopica}
\end{equation}
Cada registro se evalúa en la forma del plano donde fue producido.
Los transportes inversos de \eqref{eul:eq:recuperacion-variable}
restituyen todos los estados en orden inverso a partir del terminal y
los registros completos.

Una realización concreta utiliza matrices invertibles \(M_n\),
\(G_n=M_n^{-\mathsf T}M_n^{-1}\) y rotaciones \(R_n\):
\[
 B_n=M_{n+1}M_n^{-1},\qquad
 C_n=M_{n+1}R_nM_n^{-1}.
\]
Las dos identidades de formas se verifican sustituyendo estas
expresiones y utilizando \(R_n^\mathsf TR_n=I\). Las coordenadas,
las escalas y los parámetros de \(M_n,R_n\) se reciben de los lectores
del estado cuando se aplica esta realización a Mecánica Dimensional.

\subsubsection{Criterio exacto de transferencia al registro}

Si \(R_n\) es una rotación de ángulo \(\theta_n\), en coordenadas
normalizadas \(y_n=M_n^{-1}v_n\) se cumple
\[
 y_{n+1}=\frac{8I+R_n}{9}y_n,
 \qquad
 \left(\frac{8I+R_n}{9}\right)^\mathsf T
 \left(\frac{8I+R_n}{9}\right)=r_nI,
\]
\[
 r_n=1-\frac{32}{81}\sin^2(\theta_n/2).
\]
Por inducción,
\(I_N(v_N)=I_0(v_0)\prod_{n<N}r_n\).

\begin{proposicion}[Criterio de anulación del terminal]
Si \(I_0(v_0)>0\), el terminal satisface \(I_N(v_N)\to0\) exactamente
cuando
\[
 \sum_{n\ge0}\sin^2(\theta_n/2)=+\infty.
\]
\end{proposicion}
\begin{proof}
Para \(c=32/81\) y \(u\in[0,1]\), la integración de
\(1/(1-s)\) entre cero y \(cu\) da
\[
 cu\le-\log(1-cu)\le\frac{cu}{1-c}.
\]
La suma de los logaritmos diverge exactamente cuando diverge la suma
de los \(u_n=\sin^2(\theta_n/2)\). El producto de los \(r_n\)
tiende a cero exactamente en ese caso.
\end{proof}

La conservación \eqref{eul:eq:telescopica} vale a toda profundidad
finita. El criterio anterior determina cuándo toda la forma inicial
queda transferida al registro en esta realización. La secuencia angular
es la obtenida por el lector correspondiente de la trayectoria.

\subsection{Acción integrada y compatibilidad con el refinamiento}
\label{eul:accion-refinamiento}

En un plano canónico con coordenadas \((Q,P)\), sea \(\gamma\) una curva
cerrada a trozos \(C^1\). Para cualquier matriz real \(L\) de orden dos,
\begin{equation}
 \oint_{L\gamma}P\,\mathrm dQ
 =\det(L)\oint_\gamma P\,\mathrm dQ.
\label{eul:eq:accion-determinante}
\end{equation}
Para demostrarlo, escríbanse \(Q'=aQ+bP\), \(P'=cQ+dP\).
Las integrales cerradas de \(Q\,\mathrm dQ\) y \(P\,\mathrm dP\)
son cero. La integración de \(\mathrm d(PQ)\) da
\(\oint Q\,\mathrm dP=-\oint P\,\mathrm dQ\). Queda el coeficiente
\(ad-bc\), que prueba la fórmula.

Aplicada a \eqref{eul:eq:determinantes}, resulta
\[
 \oint_\gamma P\,\mathrm dQ
 =\oint_{T\gamma}P\,\mathrm dQ
  +\oint_{D\gamma}P\,\mathrm dQ.
\]
La orientación de cada curva conserva el signo de su contribución.

El refinamiento requiere además compatibilidad entre los operadores
de niveles sucesivos. En este apartado, sean
\(\iota_n:V_n\to V_{n+1}\) inyecciones lineales y
\(A_n:V_n\to V_n\) endomorfismos que satisfagan
\(A_{n+1}\iota_n=\iota_nA_n\). Definimos en cada nivel
\(T_n=(8I_{V_n}+A_n)/9\) y
\(D_n=\sqrt8(A_n-I_{V_n})/9\). De estas expresiones se deduce
\[
 T_{n+1}\iota_n=\iota_nT_n,
 \qquad D_{n+1}\iota_n=\iota_nD_n.
\]
La restitución \eqref{eul:eq:recuperacion} conmuta entonces con el
refinamiento. Esta identidad conserva simultáneamente la operación
y el registro; su verificación en una realización concreta utiliza
el mapa \(\iota_n\) producido por esa prolongación.

\subsubsection{Conservación del factor escalar del lector}

Si una realización recibe un factor de escala \(e^{-x}\), el transporte
es \(C=e^{-x}E_\tau(t)\). Su determinante vale \(e^{-2x}\), y la misma
expansión de la ley nonádica da
\[
 T^\mathsf T\Omega T+D^\mathsf T\Omega D
 =\frac{8+e^{-2x}}9\Omega.
\]
Para \(x\ge0\), el complemento de escala
\[
 D_{\rm esc}=\frac{\sqrt{1-e^{-2x}}}{3}I
\]
restituye el balance completo:
\[
 \Omega=T^\mathsf T\Omega T+D^\mathsf T\Omega D
       +D_{\rm esc}^\mathsf T\Omega D_{\rm esc}.
\]
El coeficiente se obtiene del defecto determinante de la escala. Cada
aplicación conserva el factor de su lector y especifica si necesita
este registro adicional.

La extensión triádica de Euler queda así articulada con operaciones
generadas, formas invariantes y recuperadores. El ciclo APP, el
transporte nilpotente y la incidencia modal conservan su procedencia;
la realización analítica describe sus evoluciones; la división
nonádica de la lectura conserva el complemento necesario para la
restitución. Esta secuencia expresa el vínculo entre régimen,
transformación y memoria sin sustituir la genealogía por una fórmula
exponencial aislada.
